\section{BKP hierarchy}\label{section4}

Let $\tau(t)$ be the $\tau$-function of the BKP hierarchy.
In this case, the time variable is $t=(t_1,t_3,t_5,\dots)$.
We set
\begin{gather*}
[\alpha]_{\rm o}=\left(\alpha,\frac{\alpha^3}{3},\frac{\alpha^5}{5},\dots\right),\hskip1cm\tilde{\xi}
(t,\lambda)=\sum_{n=1}^{\infty}t_{2n-1}\lambda^{2n-1}.
\end{gather*}
The BKP hierarchy \cite{DJKM1,JM1} is def\/ined by
\begin{gather}
\oint e^{\tilde{\xi}(t-t',\lambda)}\tau\big(t-2\big[\lambda^{-1}\big]_{\rm o}\big)\tau\big(t'+2\big[\lambda^{-1}\big]_{\rm o}\big)\frac{d\lambda}
{2\pi i\lambda}=\tau(t)\tau(t').
\label{eq:50}
\end{gather}
Set $t=x-y$, $t'=x+y$.
We get
\begin{gather}
\oint e^{-2\tilde{\xi}(y,\lambda)}\tau\big(x-y-2\big[\lambda^{-1}\big]_{\rm o}\big)\tau\big(x+y+2\big[\lambda^{-1}\big]_{\rm o}\big)\frac{d\lambda}
{2\pi i\lambda}=\tau(x-y)\tau(x+y).
\label{eq:51}
\end{gather}
Set
\begin{gather*}
y=\sum_{i=1}^{n} [\alpha_i]_{\rm o}
\end{gather*}
in~\eqref{eq:51}. Then we have
\begin{gather*}%\label{eq:52}
\oint e^{-2\tilde{\xi}\big(\sum\limits_{i=1}^{n} [\alpha_i]_{\rm o},\lambda\big)}\tau\left(x-\sum_{i=1}^{n}
 [\alpha_i]_{\rm o}-2\big[\lambda^{-1}\big]_{\rm o}\right)
\tau\left(x+\sum_{i=1}^{n} [\alpha_i]_{\rm o}+2\big[\lambda^{-1}\big]_{\rm o}\right)\frac{d\lambda}{2\pi i\lambda}
\\
\qquad
{}=\tau\left(x-\sum_{i=1}^{n} [\alpha_i]_{\rm o}\right)\tau\left(x+\sum_{i=1}^{n} [\alpha_i]_{\rm o}\right).
\end{gather*}
By decomposing $-2\sum\limits_{n=1}^\infty t_{2n-1}{\lambda}^{2n-1}$ as
\begin{gather*}
-2\sum_{n=1}^\infty t_{2n-1}{\lambda}^{2n-1}=-\sum_{n=1}^{\infty}t_n{\lambda}^n+\sum_{n=1}^{\infty}
t_n(-\lambda)^n,
\end{gather*}
we get
\begin{gather*}
\exp\left(-2\tilde{\xi}\left(\sum_{i=1}^{n} [\alpha_i]_{\rm o},\lambda\right)\right)=\prod_{i=1}^{n}
\frac{1-\alpha_i\lambda}{1+\alpha_i\lambda}.
\end{gather*}
Computing the integral by taking residues as before, shifting $x$ as $x+\sum\limits_{i=1}^{n}[\alpha_i]_{\rm o}$ and
divi\-ding by~$\tau(x)^2$ we have
\begin{gather}
\sum_{i=1}^{n}(-1)^{i-1}\frac{\tau(x+2[\alpha_i]_{\rm o})}{\tau(x)}A_{1\dots\hat{i}\dots n}^{-1}
\frac{\tau\left(x+2\sum\limits_{l=1,\, l\neq i}^n[\alpha_l]_{\rm o}\right)}{\tau(x)}\nonumber
\\
\qquad
{}-A_{1\dots n}^{-1}\frac{\tau\left(x+2\sum\limits_{l=1}^{n}[\alpha_l]_{\rm o}\right)}{\tau(x)}=0,\qquad  n \ \text{odd},
\label{eq:53}
\\
\sum_{i=1}^{n-1}(-1)^{i-1}\frac{\alpha_{i,n}}{\tilde{\alpha}_{i,n}}
\frac{\tau(x+2[\alpha_i]_{\rm o}+2[\alpha_n]_{\rm o})}{\tau(x)}A_{1\dots\hat{i}\dots n-1}^{-1}
\frac{\tau\left(x+2\sum\limits_{l=1,\, l\neq i}^{n-1}[\alpha_l]_{\rm o}\right)}{\tau(x)}\nonumber
\\
\qquad
{} -A_{1\dots n}^{-1}\frac{\tau\left(x+2\sum\limits_{l=1}^{n}[\alpha_l]_{\rm o}\right)}{\tau(x)}=0,\qquad  n \ \text{even}.
\label{eq:54}
\end{gather}
Here $A_{1\dots n}$ is def\/ined by
\begin{gather*}
A_{1\dots n}=\prod_{i<j}^{n}\frac{\tilde{\alpha}_{ij}}{\alpha_{ij}},\qquad \tilde{\alpha}_{ij}
=\alpha_i+\alpha_j,\qquad \alpha_{ij}=\alpha_i-\alpha_j.
\end{gather*}
\begin{example}\label{example3}
The case $n=3$ of~\eqref{eq:53} becomes
\begin{gather}
\frac{\tau\left(x+2\sum\limits_{i=1}^{3}[\alpha_i]_{\rm o}\right)}{\tau(x)}=A_{123}\left\{\frac{\tau(x+2[\alpha_1]_{\rm o})}{\tau(x)}
\frac{\alpha_{23}}{\tilde{\alpha}_{23}}\frac{\tau(x+2[\alpha_2]_{\rm o}+2[\alpha_3]_{\rm o})}{\tau(x)}\right.
\nonumber
\\
\hskip36mm-\frac{\tau(x+2[\alpha_2]_{\rm o})}{\tau(x)}\frac{\alpha_{13}}{\tilde{\alpha}_{13}}
\frac{\tau(x+2[\alpha_1]_{\rm o}+2[\alpha_3]_{\rm o})}{\tau(x)}\nonumber
\\
\hskip36mm\left.
+\frac{\tau(x+2[\alpha_3]_{\rm o})}{\tau(x)}\frac{\alpha_{12}}{\tilde{\alpha}_{12}}
\frac{\tau(x+2[\alpha_1]_{\rm o}+2[\alpha_2]_{\rm o})}{\tau(x)}\right\}.
\label{eq:55}
\end{gather}
\end{example}
We call equation~\eqref{eq:55} `the four term equation of the BKP hierarchy'.
\begin{example}
The case of $n=4$ of~\eqref{eq:54} is
\begin{gather}
\frac{\tau\left(x+2\sum\limits_{i=1}^{4}[\alpha_i]_{\rm o}\right)}{\tau(x)}
=A_{1234}\left\{\frac{\alpha_{14}}{\tilde{\alpha}_{14}}
\frac{\tau(x+2[\alpha_1]_{\rm o}+2[\alpha_4]_{\rm o})}{\tau(x)}\frac{\alpha_{23}}{\tilde{\alpha}_{23}}
\frac{\tau(x+2[\alpha_2]_{\rm o}+2[\alpha_3]_{\rm o})}{\tau(x)}\right.
\nonumber
\\
\hphantom{\frac{\tau\left(x+2\sum\limits_{i=1}^{4}[\alpha_i]_{\rm o}\right)}{\tau(x)}=}{}
-\frac{\alpha_{24}}{\tilde{\alpha}_{24}}\frac{\tau(x+2[\alpha_2]_{\rm o}+2[\alpha_4]_{\rm o})}{\tau(x)}
\frac{\alpha_{13}}{\tilde{\alpha}_{13}}\frac{\tau(x+2[\alpha_1]_{\rm o}+2[\alpha_3]_{\rm o})}{\tau(x)}\nonumber
\\
\left.
\hphantom{\frac{\tau\left(x+2\sum\limits_{i=1}^{4}[\alpha_i]_{\rm o}\right)}{\tau(x)}=}{}
+\frac{\alpha_{34}}{\tilde{\alpha}_{34}}\frac{\tau(x+2[\alpha_3]_{\rm o}+2[\alpha_4]_{\rm o})}{\tau(x)}\frac{\alpha_{12}
}{\tilde{\alpha}_{12}}\frac{\tau(x+2[\alpha_1]_{\rm o}+2[\alpha_2]_{\rm o})}{\tau(x)}\right\}.
\label{eq:56}
\end{gather}
\end{example}

 As is proved in Proposition~\ref{proposition8}, equation~\eqref{eq:56} can be derived from equation~\eqref{eq:55}.
\begin{theorem}\label{theorem3}
The four term equation~\eqref{eq:55} is equivalent to the BKP hierarchy~\eqref{eq:50}.
\end{theorem}

We prove this theorem in a~similar way to the case of the KP hierarchy.

In order to prove the theorem, we use Pfaf\/f\/ians.
The def\/inition and notation of Pfaf\/f\/ians are reviewed in Appendix~\ref{appendixA}.

Let us def\/ine the components of Pfaf\/f\/ians by
\begin{gather*}%\label{eq:57}
(0,j)=\frac{\tau(x+2[\alpha_j]_{\rm o})}{\tau(x)},\qquad \hskip1cm(i,j)=\frac{\alpha_{ij}}{\tilde{\alpha}_{ij}}
\frac{\tau(x+2[\alpha_i]_{\rm o}+2[\alpha_j]_{\rm o})}{\tau(x)}.
\end{gather*}
Then it is possible to rewrite~\eqref{eq:55} and~\eqref{eq:56} as
\begin{gather}
\frac{\tau\left(x+2\sum\limits_{i=1}^{3}[\alpha_i]_{\rm o}\right)}{\tau(x)}=A_{123}(0,1,2,3),
\label{eq:58}
\\
\frac{\tau\left(x+2\sum\limits_{i=1}^{4}[\alpha_i]_{\rm o}\right)}{\tau(x)}=A_{1234}(1,2,3,4),
\label{eq:59}
\end{gather}
respectively.

The following proposition can be proved in a~similar manner to Proposition~\ref{proposition1}.
\begin{proposition}\label{proposition7}
The BKP hierarchy~\eqref{eq:50} is equivalent to~\eqref{eq:53} and~\eqref{eq:54}.
\end{proposition}
\begin{proposition}\label{proposition8}
The following equations are implied by~\eqref{eq:55}:
\begin{gather}
\frac{\tau\left(x+2\sum\limits_{i=1}^{n}[\alpha_i]_{\rm o}\right)}{\tau(x)}=A_{1\dots n}(0,1,2,\dots,n),\qquad n\ge3, \ \text{odd},
\label{eq:66}
\\
\frac{\tau\left(x+2\sum\limits_{i=1}^{n}[\alpha_i]_{\rm o}\right)}{\tau(x)}=A_{1\dots n}(1,2,\dots,n),\qquad n\ge4, \ \text{even}.
\label{eq:67}
\end{gather}
\end{proposition}

\begin{proof}
First we prove that~\eqref{eq:55} implies~\eqref{eq:56}.
Shift $x$ in~\eqref{eq:55} as $x\rightarrow x+2[\alpha_4]_{\rm o}$:
\begin{gather}
\frac{\tau\left(x+2\sum\limits_{i=1}^{4}[\alpha_i]_{\rm o}\right)}{\tau(x+2[\alpha_4]_{\rm o})}\nonumber\\
\qquad{}
=A_{123}\frac{\tau(x)^2}{\tau(x+2[\alpha_4]_{\rm o})^2}\left\{\frac{\tau(x+2[\alpha_1]_{\rm o}+2[\alpha_4]_{\rm o})}{\tau(x)}
\frac{\alpha_{23}}{\tilde{\alpha}_{23}}\frac{\tau(x+2[\alpha_2]_{\rm o}+2[\alpha_3]_{\rm o}+2[\alpha_4]_{\rm o})}{\tau(x)}
\right.
\nonumber
\\
\qquad\quad{} -\frac{\tau(x+2[\alpha_2]_{\rm o}+2[\alpha_4]_{\rm o})}{\tau(x)}\frac{\alpha_{13}}{\tilde{\alpha}_{13}}
\frac{\tau(x+2[\alpha_1]_{\rm o}+2[\alpha_3]_{\rm o}+2[\alpha_4]_{\rm o})}{\tau(x)}\nonumber
\\
 \left.
 \qquad\quad{} +\frac{\tau(x+2[\alpha_3]_{\rm o}+2[\alpha_4]_{\rm o})}{\tau(x)}\frac{\alpha_{12}}{\tilde{\alpha}_{12}}
\frac{\tau(x+2[\alpha_1]_{\rm o}+2[\alpha_2]_{\rm o}+2[\alpha_4]_{\rm o})}{\tau(x)}\right\}.
\label{eq:68}
\end{gather}
Use~\eqref{eq:58} to rewrite $\tau(x+2[\alpha_{i_1}]_{\rm o}+2[\alpha_{i_2}]_{\rm o}+2[\alpha_{i_3}]_{\rm o})$
in~\eqref{eq:68}.
Then we get~\eqref{eq:59}.
We prove~\eqref{eq:66} by induction on $n$.
The case $n=3$ is obvious.
Suppose that~\eqref{eq:66} holds in case of $n$:
\begin{gather}
\frac{\tau\left(x+2\sum\limits_{i=1}^{n}[\alpha_i]_{\rm o}\right)}{\tau(x)}=A_{1\dots n}(0,1,2,\dots,n)=A_{1\dots n}{\rm Pf}\, {\mathcal A},
\label{eq:69}
\end{gather}
where ${\mathcal A}=(a_{ij})_{0\le i,j\le n}$ is a~skew-symmetric matrix,
\begin{gather*}%\label{eq:70}
a_{ij}=
\begin{cases}\dfrac{\tau(x+2[\alpha_j]_{\rm o})}{\tau(x)}, &  i=0,
\vspace{2mm}\\
\dfrac{\alpha_{ij}}{{\tilde{\alpha}}_{ij}}\dfrac{\tau(x+2[\alpha_i]_{\rm o}+2[\alpha_j]_{\rm o})}{\tau(x)}
,&  i\neq0,\quad  i<j.
\end{cases}
\end{gather*}
In~\eqref{eq:69} we shift $x$ as
\begin{gather*}
x\rightarrow x+2[\alpha_{n+1}]_{\rm o}+2[\alpha_{n+2}]_{\rm o}.
\end{gather*}
Then we have, using~\eqref{eq:58} and~\eqref{eq:59},
\begin{gather}
\frac{\tau\left(x+2\sum\limits_{i=1}^{n+2}[\alpha_i]_{\rm o}\right)}{\tau(x)}=A_{1\dots n+2}(n+1,n+2)^{-\frac{n-1}{2}}{\rm Pf}\,
{\mathcal B},
\label{eq:71}
\end{gather}
where ${\mathcal B}=(b_{ij})_{0\le i<j\le n}$ is a~skew-symmetric matrix given by
\begin{gather*}%\label{eq:72}
b_{ij}=(n+1,n+2,i,j),\qquad  i<j .
\end{gather*}

By the analogue of the Sylvester' theorem for Pfaf\/f\/ians (Appendix~\ref{appendixB}), we have
\begin{gather*}%\label{eq:73}
{\rm Pf}((1,2,\dots,2r,i,j))_{2r+1\le i<j\le2m}=(1,2,\dots,2r)^{m-r-1}(1,2,\dots,2m).
\end{gather*}
Consider the case $r=1$ and $2m=n+3$ of this formula:
\begin{gather}
{\rm Pf}((n+1,n+2,i,j))_{0\le i<j\le n}=(n+1,n+2)^{\frac{n-1}{2}}(n+1,n+2,0,\dots,n).
\label{eq:80}
\end{gather}
Substituting~\eqref{eq:80} to~\eqref{eq:71}, we get
\begin{gather*}%\label{eq:74}
\frac{\tau\left(x+2\sum\limits_{i=1}^{n+2}[\alpha_i]_{\rm o}\right)}{\tau(x)}=A_{1\dots n+2}(0,1,\dots,n+2).
\end{gather*}
The case of even $n$ is similarly proved.
\end{proof}

\begin{proposition}\label{proposition9}
The Pl\"{u}cker's relations for Pfaffians of the right hand side of~\eqref{eq:66} and~\eqref{eq:67}
give the addition formulae~\eqref{eq:53} and~\eqref{eq:54} respectively.
\end{proposition}

\begin{proof}
Using the Pl\"{u}cker's relations~\eqref{eq:500} and~\eqref{eq:501} for Pfaf\/f\/ians given in Appendix~\ref{appendixC},
the proposition can easily be checked by direct calculations.
\end{proof}
